% ============================================================
\section{Tactile-based Grasp and Transport Anomaly Detection}
\label{sec:anomaly}
% ============================================================

\subsection{Phase-dependent anomaly definition}
\label{sec:anomaly_definition}

The tactile response of a deformable garment depends strongly on the
manipulation phase. Before gripper closure, the absence of contact is an
expected state. After closure, however, persistent absence of contact indicates
an empty or insufficient grasp. During transport, a failure may instead appear
as gradual redistribution of the contact forces, partial slip, or abrupt
contact loss. We therefore activate tactile monitoring only after gripper
closure and use separate detector configurations for the grasp-confirmation and
transport phases.

Both configurations follow a one-class formulation. The recurrent models are
trained only on successful trials and learn the tactile dynamics associated
with a stable grasp. Failure trials are excluded from model fitting and are
used only for held-out evaluation. At run time, an anomaly detected in either
phase interrupts the current motion and invokes the recovery procedure
described in Section~\ref{sec:overview}.

\subsection{Tactile preprocessing}
\label{sec:tactile_preprocessing}

The deployed detector uses one $4\times4$ three-axis tactile array. Let
$\mathbf{s}_i(t)=[s_i^x(t),s_i^y(t),s_i^z(t)]^\top$ denote the raw output of
taxel $i\in\{1,\ldots,N\}$, where $N=16$. Each sensor axis is normalized using
the calibration minimum and range:
%
\begin{equation}
    \mathbf{q}_i(t)
    =
    \frac{\mathbf{s}_i(t)-\mathbf{s}_{i,\min}}
         {\mathbf{s}_{i,\mathrm{range}}},
    \label{eq:taxel_normalization}
\end{equation}
%
where the division is performed elementwise. At the beginning of each
manipulation phase, the cumulative relative signal is reset and subsequently
updated as
%
\begin{equation}
    \mathbf{d}_i(t)
    =
    \sum_{k=t_s+1}^{t}
    \bigl(\mathbf{q}_i(k)-\mathbf{q}_i(k-1)\bigr),
    \label{eq:cumulative_tactile}
\end{equation}
%
where $t_s$ is the phase-start frame. In the absence of a sensor reset,
Eq.~\eqref{eq:cumulative_tactile} is equivalent to
$\mathbf{q}_i(t)-\mathbf{q}_i(t_s)$.

For a set of taxel vectors $\{\mathbf{v}_i(t)\}_{i=1}^{N}$, we define
$r_i(t)=\|\mathbf{v}_i(t)\|_2$ and retain taxels satisfying
$r_i(t)>\theta_c$, with $\theta_c=0.1$ in the normalized sensor space. The mean
unit direction and its mean resultant length (MRL) are
%
\begin{align}
    \bar{\mathbf{u}}(t)
    &=
    \frac{1}{|\mathcal{V}(t)|}
    \sum_{i\in\mathcal{V}(t)}
    \frac{\mathbf{v}_i(t)}{r_i(t)},                                     
    \label{eq:mean_unit_direction}\\
    \operatorname{MRL}(t)
    &=\|\bar{\mathbf{u}}(t)\|_2,
    \label{eq:implemented_mrl}
\end{align}
%
where $\mathcal{V}(t)=\{i:r_i(t)>\theta_c\}$. We set both quantities to zero
when fewer than two taxels are active. This unweighted directional statistic
prevents large-magnitude taxels from dominating the spatial-coherence
measurement.

\subsection{Grasp-confirmation detector}
\label{sec:grasp_detector}

\subsubsection{Channel A: 6D GRU autoencoder}

During grasp confirmation, Channel A receives the following six-dimensional
feature vector computed from $\mathbf{d}_i(t)$:
%
\begin{equation}
    \mathbf{x}^{G}(t)
    =
    \left[
    z(t),\,
    \rho(t),\,
    r_{\max}(t),\,
    \bar{u}_x(t),\,
    \bar{u}_y(t),\,
    \bar{u}_z(t)
    \right]^\top ,
    \label{eq:grasp_6d}
\end{equation}
%
where $z(t)=\operatorname{MRL}(t)$,
$r_{\max}(t)=\max_i\|\mathbf{d}_i(t)\|_2$, and
%
\begin{equation}
    \rho(t)
    =
    \left\|
    \frac{\sum_{i=1}^{N}\mathbf{d}_i(t)}
         {\sum_{i=1}^{N}\|\mathbf{d}_i(t)\|_2+\epsilon}
    \right\|_2 .
    \label{eq:grasp_resultant_ratio}
\end{equation}
%
The features are standardized using the mean and standard deviation estimated
from the successful training trials.

The grasp model is a two-layer GRU autoencoder with hidden dimension 128 and
dropout 0.2. Given a window
$\mathbf{X}^{G}_t\in\mathbb{R}^{T\times6}$ of length $T=50$, the encoder maps
the sequence to its final hidden state and the decoder reconstructs all
$T$ feature vectors. The anomaly score is the mean squared reconstruction
error
%
\begin{equation}
    s_G(t)
    =
    \frac{1}{6T}
    \sum_{k=1}^{T}
    \left\|
    \mathbf{x}^{G,\mathrm{std}}_{t-T+k}
    -
    \hat{\mathbf{x}}^{G,\mathrm{std}}_{t-T+k}
    \right\|_2^2 .
    \label{eq:grasp_score}
\end{equation}
%
Inference is evaluated every $q_G=10$ accepted tactile frames after the first
window has been filled. The reconstruction threshold $\tau_G$ is calibrated
from reconstruction errors of a stable-grasp calibration set. The trial-level
separation of the training, calibration, and test sets is specified in
Section~\ref{sec:experiments}.

The deployed grasp logic additionally contains two guard conditions. First, an
empty-grasp guard is activated when the maximum MRL observed after closure
remains below $\tau_{\mathrm{MRL}}=0.3$. Second, a reconstruction score greater
than $\kappa\tau_G$, with $\kappa=2.2$, is treated as a severe anomaly without
waiting for repeated regular-threshold exceedances. Thus, the grasp Channel-A
decision is
%
\begin{equation}
    a_G(t)
    =
    \mathbb{I}\!\left[z_{\max}(t)<\tau_{\mathrm{MRL}}\right]
    \lor
    \mathbb{I}\!\left[s_G(t)>\kappa\tau_G\right]
    \lor
    \mathbb{I}\!\left[s_G(t)>\tau_G\right]_{\times n_G},
    \label{eq:grasp_channel_a}
\end{equation}
%
where $\mathbb{I}[\cdot]_{\times n}$ denotes persistence for $n$ consecutive
evaluations. The current grasp checkpoint uses $n_G=1$.

\subsubsection{Channel B: framewise MRL-change detector}

The framewise grasp channel monitors changes in MRL. Let
$\mathcal{H}_t=\{z(t-W),\ldots,z(t-1)\}$ be the recent history with $W=3$.
Its framewise change score is
%
\begin{equation}
    \delta_G(t)
    =
    \left|
    z(t)-\frac{1}{W}\sum_{k=1}^{W}z(t-k)
    \right|.
    \label{eq:grasp_delta}
\end{equation}
%
The first 30 frames are used as a warm-up period. To suppress unstable MRL
values under negligible force, the detector is evaluated only when
$r_{\max}(t)>0.3$. Channel B is activated when
$\delta_G(t)>\tau_{\delta G}$ for five consecutive frames. The threshold
$\tau_{\delta G}$ is calibrated from successful grasp trials.
In the current online implementation, a confirmed Channel-B flag is latched
and consumed at the next Channel-A inference update. Consequently, the
grasp-stage response time must be measured from the complete online system and
must not be equated directly with the five-frame confirmation length.

\subsection{Transport-stage detector}
\label{sec:transport_detector}

\subsubsection{Stable-grasp reference}

The cumulative signal in Eq.~\eqref{eq:cumulative_tactile} contains the
contact state present at the beginning of transport. To make the detector
sensitive to changes from the held-garment state rather than to the absolute
contact configuration, we estimate a taxel-wise stable-grasp reference from
the first $B=30$ transport frames:
%
\begin{equation}
    \mathbf{b}_i
    =
    \operatorname{median}_{1\leq t\leq B}\mathbf{d}_i(t).
    \label{eq:transport_baseline}
\end{equation}
%
The transport change vector is then
%
\begin{equation}
    \Delta\mathbf{f}_i(t)=\mathbf{d}_i(t)-\mathbf{b}_i.
    \label{eq:transport_delta_vector}
\end{equation}
%
The same fixed reference is used for all subsequent frames in the transport
episode.

\subsubsection{Channel A: relative 9D GRU autoencoder}

Using $\mathbf{v}_i(t)=\Delta\mathbf{f}_i(t)$ in
Eqs.~\eqref{eq:mean_unit_direction}--\eqref{eq:implemented_mrl}, we form the
relative nine-dimensional feature vector
%
\begin{equation}
    \mathbf{x}^{T}(t)
    =
    \left[
    z_\Delta(t),\,
    \mathbf{F}_\Delta(t)^\top,\,
    \mathbf{M}_\Delta(t)^\top,\,
    A_\Delta(t),\,
    r_{\Delta,\max}(t)
    \right]^\top ,
    \label{eq:transport_9d}
\end{equation}
%
where
%
\begin{align}
    \mathbf{F}_\Delta(t)
    &=\sum_{i=1}^{N}\Delta\mathbf{f}_i(t),                             
    \label{eq:delta_force}\\
    \mathbf{M}_\Delta(t)
    &=\sum_{i=1}^{N}\mathbf{p}_i\times\Delta\mathbf{f}_i(t),           
    \label{eq:delta_moment}\\
    A_\Delta(t)
    &=\sum_{i=1}^{N}
    \sigma\!\left(
    \alpha\frac{\|\Delta\mathbf{f}_i(t)\|_2-\theta_c}{\theta_c}
    \right),                                                          
    \label{eq:changed_taxel_count}\\
    r_{\Delta,\max}(t)
    &=\max_i\|\Delta\mathbf{f}_i(t)\|_2.                               
    \label{eq:peak_delta_norm}
\end{align}
%
Here, $\mathbf{p}_i$ is the planar position of taxel $i$ relative to the
center of the $4\times4$ array, $\sigma(\cdot)$ is the logistic sigmoid, and
$\alpha=10$. The soft count $A_\Delta$ represents the spatial extent of
contact change; it is not interpreted as an absolute physical contact area.

Each feature dimension is standardized using statistics fitted only to the
successful training windows. To reduce between-window offsets, the first
standardized feature vector is then subtracted from every frame in the
window:
%
\begin{equation}
    \tilde{\mathbf{x}}^{T}_k
    =
    \mathbf{x}^{T,\mathrm{std}}_k
    -
    \mathbf{x}^{T,\mathrm{std}}_1,
    \qquad k=1,\ldots,T.
    \label{eq:first_frame_relative}
\end{equation}
%
The transport model is a one-layer GRU autoencoder with hidden dimension 32,
window length $T=50$, and no dropout. It is trained by minimizing the
reconstruction MSE of stable transport windows. Its anomaly score is
%
\begin{equation}
    s_T(t)
    =
    \frac{1}{9T}
    \sum_{k=1}^{T}
    \left\|
    \tilde{\mathbf{x}}^{T}_{t-T+k}
    -
    \hat{\mathbf{x}}^{T}_{t-T+k}
    \right\|_2^2 .
    \label{eq:transport_score}
\end{equation}
%
Training windows are extracted with a stride of 10 frames, and online
inference also uses $q_T=10$. The regular anomaly threshold is the
$p_A$-th percentile of reconstruction scores on successful validation
windows:
%
\begin{equation}
    \tau_T
    =
    \operatorname{Percentile}_{p_A}
    \left(\{s_T^{\mathrm{val}}\}\right).
    \label{eq:transport_threshold}
\end{equation}
%
For the current checkpoint, $p_A=96$ and $\tau_T=0.126551$. A regular
exceedance must persist for $n_T=2$ consecutive inference windows, whereas
$s_T(t)>2.2\tau_T$ triggers immediately. Because this feature is already
defined relative to a stable held-grasp reference, the grasp-stage MRL
empty-contact guard is disabled during transport.

\subsubsection{Channel B: absolute contact-loss detector}
\label{sec:transport_channel_b}

Channel A models changes in the relative tactile dynamics. To directly detect
loss of absolute contact, transport Channel B uses the raw sensor level after
removing the empty-gripper offset. Before grasping, an unloaded reference
$\mathbf{s}_{i,0}$ is estimated as the median of 20 stationary samples. The
zero-corrected and range-normalized taxel vector is
%
\begin{equation}
    \mathbf{g}_i(t)
    =
    \frac{\mathbf{s}_i(t)-\mathbf{s}_{i,0}}
         {\mathbf{s}_{i,\mathrm{range}}}.
    \label{eq:zero_corrected_force}
\end{equation}
%
The contact threshold is adapted to the measured unloaded noise:
%
\begin{equation}
    \theta_0
    =
    \max\left(
    0.03,\,
    3\,\operatorname{Percentile}_{99}
    \{\|\mathbf{g}_i(t)\|_2\}_{\mathrm{unloaded}}
    \right).
    \label{eq:contact_noise_threshold}
\end{equation}
%
We compute a soft contact count and a robust force-magnitude statistic:
%
\begin{align}
    C(t)
    &=
    \sum_{i=1}^{N}
    \sigma\!\left(
    10\frac{\|\mathbf{g}_i(t)\|_2-\theta_0}{\theta_0}
    \right),                                                          
    \label{eq:absolute_contact_count}\\
    P_3(t)
    &=
    \operatorname{mean}
    \left(
    \operatorname{Top3}_{i}
    \|\mathbf{g}_i(t)\|_2
    \right).
    \label{eq:top3_force}
\end{align}
%
Their stable-grasp references, $C_0$ and $P_{3,0}$, are the medians of the
first 30 transport frames. The relative contact and force drops are
%
\begin{align}
    d_C(t)
    &=\max\left(0,\frac{C_0-C(t)}{C_0}\right),                         
    \label{eq:contact_drop}\\
    d_P(t)
    &=\max\left(0,\frac{P_{3,0}-P_3(t)}{P_{3,0}}\right).               
    \label{eq:force_drop}
\end{align}
%
Channel B is positive when $d_C(t)>0.25$ or $d_P(t)>0.30$. An anomaly is
confirmed after five consecutive positive frames. A severe drop,
$\max(d_C(t),d_P(t))>0.60$, is confirmed immediately.

\subsection{Decision fusion and online execution}
\label{sec:tactile_fusion}

Within each phase, the final tactile decision is the logical OR of its two
channels:
%
\begin{equation}
    a_{\mathrm{phase}}(t)
    =
    a_{\mathrm{phase},A}(t)
    \lor
    a_{\mathrm{phase},B}(t).
    \label{eq:phase_fusion}
\end{equation}
%
The channels are complementary in both input representation and temporal
behavior. The GRU channel detects multivariate temporal patterns that deviate
from successful grasp or transport sequences. The framewise channel provides
an additional response to abrupt changes: MRL change during grasp confirmation
and absolute contact loss during transport. During grasp confirmation its flag
is consumed at the next recurrent-model update, whereas transport Channel B
can interrupt the motion immediately after its own confirmation rule is
satisfied.

For transport, the first 30 frames establish both the relative 9D baseline and
the absolute-contact baseline. Channel A first becomes available after its
50-frame window is filled, while Channel B is evaluated on every subsequent
accepted tactile frame. Once either channel confirms an anomaly, the current
motion is halted and the recovery state machine performs re-localization and
re-grasping. Detection latency is reported in Section~\ref{sec:experiments}
using sensor timestamps and the actual motion-stop time rather than converting
frame counts using the nominal sensor publication rate.

\begin{table*}[t]
\caption{Phase-dependent configuration of the deployed tactile anomaly
detector. Thresholds marked ``validation'' are calibrated using successful
validation trials.}
\label{tab:tactile_detector_configuration}
\centering
\small
\resizebox{\textwidth}{!}{%
\begin{tabular}{lllll}
\hline
Phase and channel & Input & Temporal context & Confirmation & Detection target \\
\hline
Grasp A
& 6D tactile feature
& 50-frame GRU-AE, stride 10
& validation threshold; severe and empty-grasp guards
& empty or unstable grasp \\
Grasp B
& MRL change
& 3-frame reference
& 5 consecutive frames
& abrupt directional change \\
Transport A
& stable-baseline relative 9D
& 50-frame GRU-AE, stride 10
& $p96$, 2 consecutive windows; $2.2\tau_T$ severe cutoff
& multivariate contact-pattern deviation \\
Transport B
& contact-count and Top-3 force drops
& 30-frame stable baseline
& 5 consecutive frames; 60\% severe cutoff
& abrupt absolute contact loss \\
\hline
\end{tabular}
}
\end{table*}
